% -----------------------------
% EDITAL 1/2018 - PROGRAD/PROREC-UTFPR-FB
% Modelo de TCC Confeccionado para Edital
% conforme normas ABNT - NBR 14724
%
% Template ajustado por Francisco Reinaldo 
%		(https://orcid.org/0000-0001-6161-6755)
%		(http://lattes.cnpq.br/7401534350061823)
% -----------------------------
% New Branch: 26/jun/18 v5
% 16/11/17 versao2.7 PDF/A
% -----------------------------
% Agradecimentos a Overleaf pela oportunidade e ao grupo abntex2!
% -----------------------------

% --- LianTze Lim's Fine Tuning
\RequirePackage{scrlfile}
% \AfterClass{memoir}{\usepackage[a-3b]{pdfx}} %v.3
\BeforePackage{hyperref}{\usepackage[a-3b]{pdfx}} %v.3.2
% Ajuda da Dr. LianTze, que abntex2 não conseguiu resolver
% ---

\documentclass[
	% -- opções da classe memoir --
	12pt,				% tamanho da fonte
	openright,			% capítulos começam em pág ímpar (insere página vazia caso preciso)
	oneside,			% para impressão em frente. Oposto a twoside (frente,costas)
	a4paper,			% tamanho do papel. 
	% -- opções da classe abntex2 --
	chapter=TITLE,		% títulos de capítulos convertidos em letras maiúsculas
	%section=TITLE,		% títulos de seções convertidos em letras maiúsculas
	%subsection=TITLE,	% títulos de subseções convertidos em letras maiúsculas
	%subsubsection=TITLE,% títulos de subsubseções convertidos em letras maiúsculas
	% -- opções do pacote babel --
	english,			% idioma adicional para hifenização
	french,				% idioma adicional para hifenização
	spanish,			% idioma adicional para hifenização
	brazil				% o último idioma é o principal do documento
	]{abntex2}

% ---
% PDF/A
% ---
%
% PDF/A é um padrão ISO destinado ao arquivamento a longo prazo de documentos eletrônicos. Enfatiza a autonomia e a reprodutibilidade, bem como os metadados legíveis por máquina.
%
% A pedido da UTFPR-FB-Biblioteca, você deve converter seu PDF para PDF/A. Ao final, valide seu PDF/A em
% https://www.pdf-online.com/osa/validate.aspx
%
% PDF gerado estará ok se a mensagem for: "The document does conform to the PDF/A-3b standard."


% ---
% Pacotes básicos 
% ---

\usepackage[T1]{fontenc}		% Selecao de codigos de fonte.
\usepackage[utf8]{inputenc}		% Codificacao do documento -conversão automática dos acentos
\usepackage{lmodern}			% Usa a fonte Latin Modern			
\usepackage{lastpage}			% Usado pela Ficha catalográfica
\usepackage{indentfirst}		% Indenta o primeiro parágrafo de cada seção.
\usepackage{color}				% Controle das cores
\usepackage{graphicx}			% Inclusão de gráficos
\usepackage{microtype} 			% para melhorias de justificação
% ---
%---



% ---
% Pacotes de citações
% ---
\usepackage[brazilian,hyperpageref]{backref}	 % Paginas com as citações na bibl
\usepackage[alf]{abntex2cite}	% Citações padrão ABNT

\renewcommand{\cftsectionfont}{\normalfont} %arrumando abntex2


% --- 
% Espaçamentos entre linhas e parágrafos 
% --- 
%
% O tamanho do parágrafo é dado por:
\setlength{\parindent}{1.3cm}
%
% Controle do espaçamento entre um parágrafo e outro:
\setlength{\parskip}{0.2cm}  % tente também \onelineskip

% acronyms
\usepackage{acronym} 


% --- 
% CONFIGURAÇÕES DE PACOTES
% --- 

%\renewcommand{\chaptermark}[1]{ \markboth{#1}{}  }



        

% ---
% Pacote de Formatação de URL 
% ---
\usepackage{url} %url clicáveis
\makeatletter  \def\url@leostyle{%
  \@ifundefined{selectfont}{\def\UrlFont{\sf}}{\def\UrlFont{\small\ttfamily}}}
\makeatother \urlstyle{leo}

% ---
% Pacote Gráfico 
% ---
\usepackage{graphicx} %graphbox: Extend graphicx to improve placement of graphics 
\ifpdf
  \DeclareGraphicsExtensions{%
    .png,.PNG,%
    .pdf,.PDF,%
    .jpg,.mps,.jpeg,.jbig2,.jb2,.JPG,.JPEG,.JBIG2,.JB2}
\else
  \DeclareGraphicsExtensions{.eps}
\fi
\graphicspath{{figs/}}
%\graphicspath{{subdir1/}{subdir2/}{subdir3/}...{subdirn/}}

% ---
% Configurações do pacote backref
% Usado sem a opção hyperpageref de backref
\renewcommand{\backrefpagesname}{Citado na(s) página(s):~}
% Texto padrão antes do número das páginas
\renewcommand{\backref}{}
% Define os textos da citação
\renewcommand*{\backrefalt}[4]{
	\ifcase #1 %
		Nenhuma citação no texto.%
	\or
		Citado na página #2.%
	\else
		Citado #1 vezes nas páginas #2.%
	\fi}%
% ---

% ---
% Ajustes em tam. de fonte de Titulo1, mas gera problema em cascata
%\renewcommand{\ABNTEXchapterfont}{\bfseries\sffamily}
\renewcommand{\ABNTEXchapterfontsize}{\Large}
%\renewcommand{\ABNTEXsectionfont}{\normalfont\sffamily}


%--------------------------------------------------------------------
%--------------------------------------------------------------------
% --- 
% DADOS BASICOS DE AUTORIA
% --- 
\autor{João José Maria da Silva}
\titulo{Lógica e Aritmética na Filosofia da Matemática de Frege}
\data{2017}
\local{Francisco Beltrão, Paraná}
 
 
% --- 
% CAPA
% --- 
\renewcommand{\imprimircapa}{%
  \begin{capa}%
    \center
    \ABNTEXchapterfont\Large UNIVERSIDADE TECNOLÓGICA FEDERAL DO PARANÁ
    \par
    CAMPUS FRANCISCO BELTRÃO
    \par
    CURSO DE LICENCIATURA EM INFORMÁTICA
    
    {\vspace*{2cm} \ABNTEXchapterfont\large\imprimirautor}

    \vfill
    \begin{center}
    \ABNTEXchapterfont\bfseries\LARGE\imprimirtitulo
    \vspace*{3em}
    \end{center}
    \vfill
    
    \large\imprimirlocal

    \large\imprimirdata
    
    \vspace*{1cm}
  \end{capa}
} 
%%%%%%%%%%%%%%%
 
 
% --- 
% FOLHA DE ROSTO
% --- 
\tipotrabalho{monografia}

\preambulo{Trabalho de Conclusão de Curso, apresentado a Universidade Tecnológica Federal – Campus Francisco Beltrão, como parte das exigências para a obtenção do título de Licenciado em Informática.} 

\setlength{\ABNTEXsignwidth}{12cm} %ajuste tamanho linha textual horizontal se necessário
\setlength{\ABNTEXsignthickness}{0pt} %coloque 1pt se desejar linha. Comumente, linha de assinatura é utilizada para pessoas não letradas.
\setlength{\ABNTEXsignskip}{0.7cm} %espaçamento entre assinaturas

\orientador{Prf. Doutor. Eng. Francisco A. F. Reinaldo}
\coorientador{Prof. Dr. Apl. James Tiberius Kirk}
%%%%%%%%%%%%%%%


% ---
% compila o indice
% ---
\renewcommand{\familydefault}{\sfdefault} %suprime fonte em abntex e força nova fonte em todo o doc.
 





%-------------------------------------------------------
\begin{document}
% Seleciona o idioma do documento (conforme pacotes do babel)
\selectlanguage{brazil}

% ----------------------------------------------------------
% ELEMENTOS PRÉ-TEXTUAIS
% ----------------------------------------------------------
% \pretextual

% ---
% Capa
% ---
\imprimircapa
% ---

% ---
% Folha de rosto
% (o * indica que haverá a ficha bibliográfica)
% ---
\imprimirfolhaderosto*
% ---



% ---
% RESUMOS
% ---

% resumo em português
\setlength{\absparsep}{18pt} % ajusta o espaçamento dos parágrafos do resumo
\begin{resumo} 
 Nos \textit{Fundamentos da Aritmética} (§68), Frege propõe definir explicitamente o operador-abstração `o número de...' por meio de extensões e, a partir desta definição, provar o Princípio de Hume (\textbf{PH}). Contudo, a prova imaginada por Frege depende de uma fórmula (\textbf{BB}) não derivável no sistema em 1884. Acreditamos que a distinção entre sentido e referência e a introdução dos valores de verdade como objetos foram motivadas para justificar a introdução do Axioma IV, a partir do qual um análogo de (\textbf{BB}) é provável. Com (\textbf{BB}) no sistema, a prova do Princípio de Hume estaria garantida. Concomitantemente, percebemos que uma teoria unificada das extensões só é possível com a distinção entre sentido e referência e a introdução dos valores de verdade como objetos. Caso contrário, Frege teria sido obrigado a introduzir uma série de \textbf{Axiomas V} no seu sistema, o que acarretaria problemas com a identidade (Júlio César). Com base nestas considerações, além do fato de que, em 1882, Frege provara as leis básicas da aritmética (carta a Anton Marty), parece-nos perfeitamente plausível que estas provas foram executadas adicionando-se o \textbf{PH} ao sistema lógico de Begriffsschrift. Mostramos que, nas provas dos axiomas de Peano a partir de \textbf{PH} dentro da conceitografia, nenhum uso é feito de (\textbf{BB}). Destarte, não é necessária a introdução do Axioma IV no sistema e, por conseguinte, não são necessárias a distinção entre sentido e referência e a introdução dos valores de verdade como objetos. Disto, podemos concluir que, provavelmente, a introdução das extensões nos \textit{Fundamentos} foi um ato tardio; e que Frege não possuía uma prova formal de \textbf{PH} a partir da sua definição explícita. Estes fatos também explicam a demora na publicação das \textit{Leis Básicas da Aritmética} e o descarte de um manuscrito quase pronto (provavelmente, o livro mencionado na carta a Marty). 
\vspace{\onelineskip} 

\noindent \textbf{Palavras-chave}: Axioma IV. Axioma V. Princípio de Hume. Valores de Verdade. Gottlob Frege. 
\end{resumo}



% ---
% inserir lista de ilustrações
% ---
\pdfbookmark[0]{\listfigurename}{lof}
\listoffigures*
\cleardoublepage
% ---

% ---
% inserir lista de tabelas
% ---
\pdfbookmark[0]{\listtablename}{lot}
\listoftables*
\cleardoublepage
% ---

% ---
% inserir lista de abreviaturas e siglas em ordem alfabética
% ---
\begin{siglas} 
	\item[ICHS] Instituto de Ciências Humanas e Sociais 
	\item[LPM] Lógica Proposicional Modal 
	\item[LQM] Lógica Quantificacional Modal 
	\item[UFRJ] Universidade Federal do Rio de Janeiro
\end{siglas}

% ---
% inserir lista de símbolos
% ---
\begin{simbolos}
	\item[$ \Gamma $] Letra grega Gama
	\item[$ \Lambda $] Lambda
	\item[$ \zeta $] Letra grega minúscula zeta
	\item[$ \in $] Pertence
\end{simbolos}


% ---
% inserir o sumario
% ---
\pdfbookmark[0]{\contentsname}{toc}
\tableofcontents*
\cleardoublepage
% ---


% ----------------------------------------------------------
% ELEMENTOS TEXTUAIS
% ----------------------------------------------------------

% ------ Retira do cabeçalho o chapter name ------%
\makepagestyle{abntheadings}
\makeoddhead{abntheadings}{}{}{\ABNTEXfontereduzida\thepage}
% ------------------------------------------------%

\textual 

\chapter{\textbf{TEMA}} %inserir \textbf{} (Falha do abntex2)

Apresenta o assunto principal, o estado da arte e o estado da prática, em especial – observe o exemplo. Sobretudo, verifique as normas de citação. Sugere-se entre uma página a duas páginas.

\section[DELIMITAÇÃO DO TEMA]{DELIMITAÇÃO DO TEMA}

Indicar, sobretudo a limitação de escopo – propósito, intento efetivo, foco; a limitação geográfica – tipo de empresa, setor, local. A delimitação pode ser feita em aproximadamente meia página.  

Diversas ferramentas envolvendo gamificação tais como o \textit{Mathlab}{\textsuperscript{\textcopyright}, e o \textit{Khan Academy}{\textsuperscript{\textcopyright} \cite{Marchetto2016}, vêm sendo propostas para realizar uma melhora na aplicação de conteúdo em sala de aula.

\section[PROBLEMAS E PREMISSAS]{PROBLEMAS E PREMISSAS}

Descreve-se o contexto das principais dificuldades encontradas. Em geral incorpora-se a pergunta de pesquisa e a premissa ou hipótese levantadas.

\section[OBJETIVOS]{OBJETIVOS}

\subsection[Objetivo Geral]{Objetivo Geral}

Lembrar que o objetivo geral é sempre uma ação, bem delimitada, portanto exige um verbo no infinitivo.


\subsection[Objetivos Específicos]{Objetivos Específicos}

 Nos objetivos específicos usa-se apenas um verbo / ação por objetivo.

\section[JUSTIFICATIVA]{JUSTIFICATIVA}

Deve mostrar, cabalmente, a importância do estudo em questão, portanto, dispensa citações diretas ou indiretas, salvo quando da existência de dados numéricos ou históricos por exemplo.



{\vspace{5em}\let\clearpage\relax \chapter[PROCEDIMENTOS METODOLÓGICOS]{\textbf{PROCEDIMENTOS METODOLÓGICOS}}}

Classifica-se a pesquisa em pelo menos: natureza; objetivo macro; técnicas e procedimentos a serem utilizados. Tal classificação exige a utilização de referencial teórico próprio.



{\vspace{5em}\let\clearpage\relax \chapter[REFERENCIAL TEÓRICO]{	\textbf{REFERENCIAL TEÓRICO}}}

Indicam-se quais os assuntos serão tratados e principalmente qual o embasamento teórico e seus autores. Pode ser apresentada uma prévia da proposta de trabalho.


{\vspace{5em}\let\clearpage\relax \chapter[DEMONSTRATIVOS DA APLICAÇÃO E VINCULAÇÃO]{\textbf{DEMONSTRATIVOS DA APLICAÇÃO E VINCULAÇÃO}}}
\vspace{-1em}


\section[APLICAÇÃO]{DEMONSTRATIVO DA APLICAÇÃO DOS RECURSOS}

Demonstrativo da aplicação dos recursos.

\section[VINCULAÇÃO DO PROJETO]{DEMONSTRATIVO DA VINCULAÇÃO DO PROJETO COM A ÁREA DE INSERÇÃO PROFISSIONAL}

Demonstrativo da vinculação do projeto.


{\vspace{5em}\let\clearpage\relax \chapter[CRONOGRAMA]{\textbf{CRONOGRAMA}}}

Para a realização deste trabalho propõem-se o seguinte cronograma de realização das atividades:

Sugere-se inserir estes pontos: em Etapa mês
\begin{itemize}
\item Revisão de literatura
\item Definição dos procedimentos
\item Coleta de dados
\item Análise dos resultados
\item Redação da monografia
\item Correção e complementação da monografia
\item  Defesa
\end{itemize}

plausível que estas provas foram executadas adicionando-se o \textbf{PH} ao sistema lógico de Begriffsschrift. Mostramos que, nas provas dos axiomas de Peano a partir de \textbf{PH} dentro da conceitografia, nenhum uso é feito de (\textbf{BB}). Destarte, não é necessária a introdução.

\chapter{Alguns exemplos de tabelas e alienas}

\begin{figure}[!ht]
\centering
\caption{Logomarca ABNT}
\includegraphics[width=0.47\textwidth]{Screen_Shot_2018-06-22_at_10_14_05.png}
\fonte{sua fonte a usar}
\label{fig: Logo abnt}
\end{figure}


\begin{alineas}
\item linha 1:
\begin{alineas}
\item subalinea 1;
\item subalinea 2;
\end{alineas}
\item linha 2:
\begin{subalineas}
\item subalinea 1;
\item subalinea 2;
\end{subalineas}
\item linha 3:
\begin{incisos}
\item subalinea 1;
\item subalinea 2;
\end{incisos}
\item linha 4.
\end{alineas}


\begin{table}[htb]
\IBGEtab{%
\caption{Um Exemplo de tabela alinhada que pode ser longa ou curta,
conforme padrão IBGE.}%
\label{tabela-ibge}
}{%
\begin{tabular}{ccc}
\toprule
Nome & Nascimento & Documento \\
\midrule \midrule
Maria da Silva & 11/11/1111 & 111.111.111-11 \\
\bottomrule
\end{tabular}%
}{%
\fonte{Produzido pelos autores}%
\nota{Esta éuma nota, que diz que os dados são baseados na
regressão linear.}%
\nota[Anotações]{Uma anotação adicional, seguida de várias outras.}%
}
\end{table}

\begin{table}[htb]
\ABNTEXfontereduzida
\caption[Níveis de investigação]{Níveis de investigação.}
\label{tab-nivinv}
\begin{tabular}{p{2.6cm}|p{6.0cm}|p{2.25cm}|p{3.40cm}}
  %\hline
   \textbf{Nível de Investigação} & \textbf{Insumos}  & \textbf{Sistemas de Investigação}  & \textbf{Produtos}  \\
    \hline
    Meta-nível & Filosofia\index{filosofia} da Ciência  & Epistemologia &
    Paradigma  \\
    \hline
    Nível do objeto & Paradigmas do metanível e evidências do nível inferior &
    Ciência  & Teorias e modelos \\
    \hline
    Nível inferior & Modelos e métodos do nível do objeto e problemas do nível inferior & Prática & Solução de problemas  \\
   % \hline
\end{tabular}
\legend{Fonte: \citeonline{van86}}
\end{table}

% ----------------------------------------------------------
% Finaliza a parte no bookmark do PDF
% para que se inicie o bookmark na raiz
% e adiciona espaço de parte no Sumário
% ----------------------------------------------------------
\phantompart


% ----------------------------------------------------------
% ELEMENTOS PÓS-TEXTUAIS
% ----------------------------------------------------------
\postextual
% ----------------------------------------------------------

% ----------------------------------------------------------
% Referências bibliográficas
% ----------------------------------------------------------
%\bibliography{acervo.bib}

% ----------------------------------------------------------
% Glossário
% ----------------------------------------------------------
%
% Consulte o manual da classe abntex2 para orientações sobre o glossário.
%
%\glossary

% ----------------------------------------------------------
% Apêndices
% ----------------------------------------------------------

% ---
% Inicia os apêndices
% ---
\begin{apendicesenv}

% Imprime uma página indicando o início dos apêndices
\partapendices

% ----------------------------------------------------------
\chapter{apêndice 1}
% ----------------------------------------------------------

texto texto

% ----------------------------------------------------------
\chapter{apêndice 2}
% ----------------------------------------------------------

texto texto

\end{apendicesenv}
% ---


% ----------------------------------------------------------
% Anexos
% ----------------------------------------------------------

% ---
% Inicia os anexos
% ---
\begin{anexosenv}

% Imprime uma página indicando o início dos anexos
\partanexos

% ---
\chapter[CARTA DE ANUÊNCIA]{CARTA DE ANUÊNCIA DA COORDENAÇÃO DO CURSO PARA O EDITAL 1/2018 - PROGRAD/PROREC}

\vspace{3em}

{\linespread{2}
A coordenação do curso de Licenciatura em Informática  declara que a proposta de TRABALHO DE CONCLUSÃO DE CURSO 1, do aluno, abaixo relacionado, está de
acordo com os procedimentos e normas estabelecidas no curso e pode ser submetida ao Edital 1/2018 - PROGRAD/PROREC, de Apoio à Execução de Trabalhos de Conclusão de Curso, para o primeiro semestre de 2018.
}

\vspace{2em}
Título: \textbf{\imprimirtitulo}

\vspace{2em}
Aluno: \textbf{\imprimirautor }

\vspace{\fill}
\hspace{\fill} Francisco Beltrão-PR, \today.

\vspace{2em}
\begin{center}
\small 
\textbf{Profa. Doutora Maici Duarte Leite\\
\footnotesize{Coordenadora do Curso de Licenciatura em Informática}} \vspace{3em}

\textbf{\imprimirorientador\\ \footnotesize{Prof. Orientador TCC}}
\end{center}

\end{anexosenv}


\pagebreak
\bibliography{acervo.bib}
\end{document}

enumerate
\usepackage[shortlabels]{enumerate}
.
.
.
\begin{enumerate}[(a)] % (a), (b), (c), ...
\item
\end{enumerate}
.
.
.
\begin{enumerate}[a)] % a), b), c), ...
\item
\end{enumerate}